\pagestyle{fancy}
\fancyhf{}
\lhead{\textsl{\footnotesize{\nouppercase{}}}}
\rhead{\textsl{\footnotesize{\nouppercase{Sitography}}}}
\lfoot{\textsl{\footnotesize Atef JRIBI}}
\cfoot{\thepage}
\rfoot{\textsl{\footnotesize }}
\normalsize
\begin{thebibliography}{9}

\bibitem{elc0} %ch1
\texttt{\url{https://codeburst.io/so-wtf-is-git-fa7daa0e0271}}

\bibitem{elc1} %ch1
\texttt{\url{https://www.freecodecamp.org/news/what-is-html-definition-and-meaning/}}

\bibitem{elc2} %ch1
\texttt{\url{https://www.techopedia.com/definition/28243/cascading-style-sheets-level-3-css3}}


\bibitem{elc3} %ch4
\texttt{\url{https://reactjs.org/docs/getting-start}}
\bibitem{elc4} %ch1
\texttt{\url{https://www.djangoproject.com/}}
\bibitem{elc5} %ch1
\texttt{ \url{https://www.django-rest-framework.org/}}
\bibitem{elc6} %ch1
\texttt{\url{https://www.mongodb.com/atlas/database}}
\bibitem{elc7} %ch1
\texttt{\url{https://pypi.org/project/djongo/}}

\bibitem{elc8} %ch1
\texttt{\url{https://code.visualstudio.com/docs}}
\bibitem{elc9} %ch1
\texttt{\url{https://www.postman.com/}}

\bibitem{elc10} %ch1
\texttt{\url{shorturl.at/xHIU5}}

\bibitem{elc11} %ch1
\texttt{\url{https://git-scm.com/}}

\bibitem{elc12} %ch1
\texttt{\url{https://github.com/}}




\end{thebibliography}
\newpage
\pagestyle{fancy}
\fancyhf{}
\lhead{\textsl{\footnotesize{\nouppercase{\rightmark}}}}
\rhead{\textsl{\footnotesize{\nouppercase{\leftmark}}}}
\lfoot{\textsl{\footnotesize Atef JRIBI}}
\cfoot{\thepage}
\rfoot{\textsl{\footnotesize }}

